\section{Temporal prestress and the kinetic limit}
\label{sec:temporal}

Before analyzing the wave and gauge sectors we settle the status of the timelike link,
because it fixes the meaning of $\eta_4=+1$. In the earlier single-$\alpha$, $r_t\approx 0$
formulation the temporal term $\tfrac12\kappa_t L_{n4}^2$ is an isotropic kinetic energy, so
$\eta_4=+1$ looks like a mere kinetic sign rather than a genuine prestress. The refined form
$r_4=\alpha_t a$ with $0<\alpha_t<1$ resolves this.

\paragraph{The temporal link is the kinetic term.} Expanding a temporal link about the
vacuum, $Q_{n4}=a\hat e_4+\Delta_4\xi$ with $L_{n4}=a+\Delta_4\xi^4+\tfrac1{2a}\sum_i(\Delta_4\xi^i)^2+\dots$,
the quadratic part of $\tfrac12\kappa_t(L_{n4}-\alpha_t a)^2$ is
\begin{equation}
  S^{(2)}_{\text{time}}
  = \tfrac12\kappa_t\sum_n\Bigl[(\Delta_4\xi^4)^2+(1-\alpha_t)\textstyle\sum_i(\Delta_4\xi^i)^2\Bigr].
  \label{eq:s2time}
\end{equation}
In the long-wavelength (continuum) limit $\Delta_4\xi\to a\,\partial_t\xi$, this is a genuine
kinetic energy
\begin{equation}
  T=\tfrac12\kappa_t a^2\Bigl[(\partial_t\xi^4)^2+(1-\alpha_t)\textstyle\sum_i(\partial_t\xi^i)^2\Bigr],
  \label{eq:kinetic}
\end{equation}
with \emph{emergent inertia} $m=\kappa_t a^2$ derived from the temporal link stiffness---no
separately postulated mass. The discrete-to-continuum limit is the long-wavelength limit at
fixed lattice spacing, $2(1-\cos\omega)\to\omega^2$, verified numerically to converge as
$\omega\to 0$ \citep{repo}.

\paragraph{The $(1-\alpha_t)$ anisotropy is the fingerprint of genuine prestress.} The
longitudinal ($\xi^4$) temporal inertia is $\kappa_t$, independent of $r_4$; the transverse
($\xi^i$) temporal inertia is $\kappa_t(1-\alpha_t)=\rho_4/a$, i.e.\ it \emph{is} the
prestress. At $r_4=0$ ($\alpha_t\to 0$) the temporal term is isotropic
($\tfrac12\kappa_t|\Delta_4\xi|^2$)---indistinguishable from a plain inertia, the
kinetic-sign artifact. For $0<\alpha_t<1$ it is anisotropic, which only a spring under
tension can produce; a pure inertia cannot. The ratio transverse/longitudinal $=(1-\alpha_t)$
is verified exactly against the link Hessian \citep{repo}. The limits $\alpha_t\to 0$
(artifact) and $\alpha_t\to 1$ (vanishing transverse temporal stiffness, over-decoupling)
are both excluded; the working regime is $\alpha_t=1-\varepsilon_t$.

\paragraph{$\eta_4=+1$ makes $S=T-V$ and hence wave dynamics.} With $\eta_4=+1$ and
$\eta_i=-1$, the action is $S=T-V$: the temporal prestress term supplies the kinetic energy
$T$ and the spacelike prestress terms the potential $V$. The 4D-block stationarity is then
Hamilton's principle, and the Euler--Lagrange stencil is the symplectic Störmer--Verlet
integrator \citep{MarsdenWest2001}. A linearized time march reproduces waves propagating at
the transverse speed $c_T$ of Section~\ref{sec:metric} (measured $0.910$ vs.\ predicted
$0.913$ in dimensionless units), whereas flipping to the all-minus (Euclidean) sign turns
the update elliptic and unstable (amplitudes diverge), confirming that $\eta_4=+1$---opposite
to the spacelike sign---is what converts the static block into causal time evolution
\citep{repo}. The same $(1-\alpha_t)$ factor set here controls the amplitude-mode speed and
the sector split in Sections~\ref{sec:metric}--\ref{sec:split}.
